% GENERATED FILE. Edit canonical lessons, then run npm run generate:books.
% canonical-source-hash: fnv1a64:16476c855071ce2c
% canonical-lessons: FR-C371-annuler, FR-C371-reporter, FR-C371-provisoire, FR-C371-definitif, FR-C371-programmer

\chapter{Changement de programme}
\label{ch:fr-changement-de-programme}
\hlchaptermodality{371}

\begin{chapteropening}
\textbf{By the end of this chapter:} \emph{I can change a plan: cancel it, put it off, say it is temporary or final, and schedule something new.}

Complete the last lesson of chapter 371: I can change a plan: cancel it, put it off, say it is temporary or final, and schedule something new.
\end{chapteropening}

\section[annuler]{annuler — to cancel}

\label{lesson:FR-C371-annuler}



\begin{quote}
Before the new one: say the French for to adapt, then the French for supple.
\end{quote}

\subsection*{You'll want to know: annuler}
\textbf{annuler} — \textquotedblleft{}to cancel\textquotedblright{}.

\begin{quote}
Le concert est annulé à cause de la pluie.
\end{quote}

\begin{quote}
\emph{The concert is cancelled because of the rain.}
\end{quote}

\emph{Annuler une réservation} — to cancel a booking.

\begin{cousinweb}
\emph{annuler} comes from Late Latin \emph{annūllāre} \textquotedblleft{}to bring to nothing\textquotedblright{} (EtymDB), built on \emph{nūllus} \textquotedblleft{}none\textquotedblright{} (Lewis \& Short).
\end{cousinweb}

\begin{grammarlens}[title={What you've built}]
One more word for telling, explaining and describing at B1.
\end{grammarlens}

\subsection*{Your turn}
\begin{itemize}
\raggedright
  \item \emph{Say it:} \emph{annuler}
  \item \emph{Say it:} \emph{annuler}, once more
  \item \emph{Say it:} a sentence of your own with \emph{annuler}
  \item \emph{Recall:} say the French for to adapt, then the French for supple, then say \emph{annuler} again
\end{itemize}

\begin{tcolorbox}[breakable,colback=teal!4,colframe=teal!35!black,title={Before you move on}]
What does annuler mean? (\textquotedblleft{}To cancel\textquotedblright{}.) Say it once more.
\end{tcolorbox}
\section[reporter]{reporter — to postpone, to put off}

\label{lesson:FR-C371-reporter}



\begin{quote}
Before the new one: say the French for supple, then the French for to cancel.
\end{quote}

\subsection*{You'll want to know: reporter}
\textbf{reporter} — \textquotedblleft{}to postpone, to put off\textquotedblright{}.

\begin{quote}
Nous avons reporté la réunion à jeudi.
\end{quote}

\begin{quote}
\emph{We postponed the meeting to Thursday.}
\end{quote}

\emph{Reporter à plus tard} — to put off till later. The noun \emph{un reporter} is a news reporter.

\begin{cousinweb}
\emph{reporter} is built from \emph{re-} \textquotedblleft{}back\textquotedblright{} and \emph{porter} \textquotedblleft{}to carry\textquotedblright{}, after Latin \emph{reportāre} \textquotedblleft{}to carry back\textquotedblright{} (Lewis \& Short), on \emph{portāre} \textquotedblleft{}to carry\textquotedblright{}: you carry the date further on.
\end{cousinweb}

\begin{grammarlens}[title={What you've built}]
One more word for telling, explaining and describing at B1.
\end{grammarlens}

\subsection*{Your turn}
\begin{itemize}
\raggedright
  \item \emph{Say it:} \emph{reporter}
  \item \emph{Say it:} \emph{reporter}, once more
  \item \emph{Say it:} a sentence of your own with \emph{reporter}
  \item \emph{Recall:} say the French for supple, then the French for to cancel, then say \emph{reporter} again
\end{itemize}

\begin{tcolorbox}[breakable,colback=teal!4,colframe=teal!35!black,title={Before you move on}]
What does reporter mean? (\textquotedblleft{}To postpone, to put off\textquotedblright{}.) Say it once more.
\end{tcolorbox}
\section[provisoire]{provisoire — temporary, provisional}

\label{lesson:FR-C371-provisoire}



\begin{quote}
Before the new one: say the French for to cancel, then the French for to postpone.
\end{quote}

\subsection*{You'll want to know: provisoire}
\textbf{provisoire} — \textquotedblleft{}temporary, provisional\textquotedblright{}.

\begin{quote}
C'est une solution provisoire, en attendant mieux.
\end{quote}

\begin{quote}
\emph{It is a temporary solution, until something better comes.}
\end{quote}

One form for both genders.

\begin{cousinweb}
\emph{provisoire} comes from Medieval Latin \emph{prōvīsōrius}, built on Latin \emph{prōvīsus}, the participle of \emph{prōvidēre} \textquotedblleft{}to foresee, to provide for\textquotedblright{} (Lewis \& Short).
\end{cousinweb}

\begin{grammarlens}[title={What you've built}]
One more word for telling, explaining and describing at B1.
\end{grammarlens}

\subsection*{Your turn}
\begin{itemize}
\raggedright
  \item \emph{Say it:} \emph{provisoire}
  \item \emph{Say it:} \emph{provisoire}, once more
  \item \emph{Say it:} a sentence of your own with \emph{provisoire}
  \item \emph{Recall:} say the French for to cancel, then the French for to postpone, then say \emph{provisoire} again
\end{itemize}

\begin{tcolorbox}[breakable,colback=teal!4,colframe=teal!35!black,title={Before you move on}]
What does provisoire mean? (\textquotedblleft{}Temporary, provisional\textquotedblright{}.) Say it once more.
\end{tcolorbox}
\section[définitif]{définitif — final, definitive}

\label{lesson:FR-C371-definitif}



\begin{quote}
Before the new one: say the French for to postpone, then the French for temporary.
\end{quote}

\subsection*{You'll want to know: définitif}
\textbf{définitif} — \textquotedblleft{}final, definitive\textquotedblright{}.

\begin{quote}
Sa décision est définitive.
\end{quote}

\begin{quote}
\emph{Her decision is final.}
\end{quote}

Feminine \emph{définitive}. \emph{En définitive} — in the end.

\begin{cousinweb}
\emph{définitif} is a learned borrowing from Latin \emph{dēfīnītīvus} (EtymDB), from \emph{dēfīnīre} \textquotedblleft{}to set limits, to define\textquotedblright{}, built on \emph{fīnis} \textquotedblleft{}end, boundary\textquotedblright{} (Lewis \& Short).
\end{cousinweb}

\begin{grammarlens}[title={What you've built}]
One more word for telling, explaining and describing at B1.
\end{grammarlens}

\subsection*{Your turn}
\begin{itemize}
\raggedright
  \item \emph{Say it:} \emph{définitif}
  \item \emph{Say it:} \emph{définitif}, once more
  \item \emph{Say it:} a sentence of your own with \emph{définitif}
  \item \emph{Recall:} say the French for to postpone, then the French for temporary, then say \emph{définitif} again
\end{itemize}

\begin{tcolorbox}[breakable,colback=teal!4,colframe=teal!35!black,title={Before you move on}]
What does définitif mean? (\textquotedblleft{}Final, definitive\textquotedblright{}.) Say it once more.
\end{tcolorbox}
\section[programmer]{programmer — to schedule, to plan; to program}

\label{lesson:FR-C371-programmer}



\begin{quote}
Before the new one: say the French for temporary, then the French for final.
\end{quote}

\subsection*{You'll want to know: programmer}
\textbf{programmer} — \textquotedblleft{}to schedule, to plan; to program\textquotedblright{}.

\begin{quote}
Nous avons programmé la fête pour samedi.
\end{quote}

\begin{quote}
\emph{We scheduled the party for Saturday.}
\end{quote}

Also to program a machine.

\begin{cousinweb}
\emph{programmer} is built on \emph{programme}, from Late Latin \emph{programma} \textquotedblleft{}a proclamation, an edict\textquotedblright{} (EtymDB; Lewis \& Short), from Greek \emph{prógramma} \textquotedblleft{}a written public notice\textquotedblright{}, built from \emph{pró} \textquotedblleft{}before\textquotedblright{} and \emph{gráphein} \textquotedblleft{}to write\textquotedblright{}.
\end{cousinweb}

\begin{grammarlens}[title={What you've built}]
One more word for telling, explaining and describing at B1.
\end{grammarlens}

\subsection*{Your turn}
\begin{itemize}
\raggedright
  \item \emph{Say it:} \emph{programmer}
  \item \emph{Say it:} \emph{programmer}, once more
  \item \emph{Say it:} a sentence of your own with \emph{programmer}
  \item \emph{Recall:} say the French for temporary, then the French for final, then say \emph{programmer} again
\end{itemize}

\begin{tcolorbox}[breakable,colback=teal!4,colframe=teal!35!black,title={Before you move on}]
What does programmer mean? (\textquotedblleft{}To schedule, to plan; to program\textquotedblright{}.) Say it once more.
\end{tcolorbox}
